\needspace{8\baselineskip}
\section{模型假设}

本文只补充题面未给定而模型必须明确的现实简化，每条假设后附其成立依据。

（1）药材内部的温度与含水率只沿半径变化，轴向与周向的梯度可以忽略。药材长 25~cm、半径 2~cm，侧面积与两端面积之比等于长度与半径之比即 $12.5$，端面换热仅占总表面的 $7.4\%$；题目要求给出的也只是“到药材中心距离”处的结果，径向是题意认定的唯一空间方向。

（2）药材视为均质连续介质，孔隙内水分的整体迁移等效为遵循 Fick 定律的扩散过程，不区分自由水与结合水，也不单独建立气相压力场。题目已直接给出水分浓度扩散系数的经验公式，即指定了本构关系的形式，这种以有效扩散系数概括内部迁移的做法在干燥建模中被广泛采用\upcite{ding2026timevar}；若另建两相多孔介质模型，则需要引入题面未提供的渗透率与毛细压力等参数\upcite{whitaker1977porous}。

（3）能量方程只含导热项与表面对流换热，不设置蒸发相变的源汇项。题目在能量一侧只给出密度、比热容、导热系数与对流换热系数四个参数，三组附录中均无蒸发潜热，故主模型不从常识补入该参数。这一简化的影响本文如实登记而不宣称可忽略：另设一支在热方程中加入潜热汇、潜热取量级参考值 $2.4\times10^6$~J/kg，算得的烘干时长比不设潜热汇的基准慢约一成，即若题目提供潜热数据，结果应向偏慢一侧修正；具体数值见第~\ref{sec:sens} 节的对照算例。

（4）表面对流传质的驱动力取药材表面含水率与烘房空气水分浓度之差。题目对药材写“水分浓度（即干基含水率）”、对烘房同样写“水分浓度”，二者单位相同且只给出一个对流传质系数，无换算系数可用；量纲上传质系数与含水率之积同扩散系数与含水率梯度之积一致，边界条件自洽。

（5）烘房激励在实测序列覆盖的四小时之后按其尾段实测平台值延拓，取空气温度 $49.9989$~°C、空气水分浓度 $0.049988$~kg/kg。题目明示烘干包含预热平衡与恒温干燥两个阶段且全程持续两到三天，实测序列恰好覆盖前一阶段，尾段增量已落在零附近；延拓值取自实测而非拟合外推，不引入题外参数。

（6）问题四中的收缩为径向仿射收缩，轴向长度保持不变，半径随时间的变化完全由实测序列给定。题目要求“根据附件 2”确定时长，即把半径序列定为已知输入；轴向不收缩这一点还得到干物质守恒的独立支持，若改设为各向同性收缩，反推的干物质密度将大幅超出经验公式给出的绝干值。
